\documentclass[11pt,a4paper]{article}

\usepackage[T1]{fontenc}
\usepackage[utf8]{inputenc}
\usepackage{lmodern}
\usepackage{geometry}
\usepackage{enumitem}
\usepackage{tabularx}
\usepackage{booktabs}
\usepackage{eurosym}
\usepackage[table]{xcolor}
\usepackage{titlesec}
\usepackage{needspace}
\usepackage{setspace}
\usepackage{array}
\usepackage{ragged2e}
\usepackage[hidelinks]{hyperref}

\geometry{top=1.2cm,bottom=1.6cm,left=1.2cm,right=1.2cm}

\definecolor{accent}{HTML}{000000}
\definecolor{cellbg}{HTML}{F2F2F2}

\newcommand{\sectionrule}[1]{\needspace{2\baselineskip}\MakeUppercase{#1}\\[-0.3em]{\color{accent}\rule{\textwidth}{0.4pt}}}
\newcommand{\subsectionformat}[1]{\needspace{1.5\baselineskip}\MakeUppercase{#1}}
\titleformat{\section}{\normalfont\bfseries\color{accent}\centering}{ }{0pt}{\sectionrule}
\titleformat{\subsection}{\normalfont\bfseries\color{accent}}{ }{0pt}{\subsectionformat}

% ---------------------------------------------------------------
% Unified vertical-rhythm system.
% One shared length (\contentgap) governs the space left after ANY
% standalone block of content, whether that block is a CV entry
% (via \entryspace, used inside \entry/\entrywidedate/\entrynarrowdate)
% or a non-entry block such as a paragraph or table (via \blockspace,
% inserted manually). Because both draw on the same length, the gap
% before every \section is identical no matter what precedes it.
% Section/subsection "before" spacing is set once here, rather than
% re-approximated with one-off \vspace calls at each occurrence.
% ---------------------------------------------------------------
\newlength{\contentgap}
\setlength{\contentgap}{1.3em}
\newcommand{\blockspace}{\vspace{\contentgap}}

\titlespacing*{\section}{0pt}{1.0em}{0.6em}
\titlespacing*{\subsection}{0pt}{1.0em}{0.3em}

\setlist[itemize]{leftmargin=1.2em,topsep=0pt,itemsep=0pt,labelsep=0.35em,parsep=0pt,partopsep=0pt}
\setlength{\parindent}{0pt}
\setlength{\parskip}{0pt}
\setlength{\tabcolsep}{3pt}
\setlength{\emergencystretch}{2em}
\hbadness=10000
\hfuzz=1pt

\newcolumntype{L}{>{\RaggedRight\arraybackslash}p{0.74\textwidth}}
\newcolumntype{R}{>{\RaggedLeft\arraybackslash}p{0.26\textwidth}}

\newcommand{\entrybulletspace}{\vspace*{-0.2em}}
\newcommand{\entrynobulletspace}{\vspace{0.3em}}
\newcommand{\entryspace}{\vspace{\contentgap}}

% Single shared left-column width for ALL entries, so title/subtitle
% always wrap at the same right edge, with place/date reserved in a
% consistent right column (sized to fit one line each).
\newlength{\entrytitlewidth}
\setlength{\entrytitlewidth}{0.70\textwidth}

\newcommand{\entry}[5]{%
  \begin{minipage}[t]{\textwidth}
  \begingroup\setlength{\tabcolsep}{0pt}%
  \begin{tabular*}{\textwidth}{@{}p{\entrytitlewidth}@{\extracolsep{\fill}}r@{}}
    \textbf{#1} & #3 \\
    \textit{#2} & \mbox{#4} \\
  \end{tabular*}
  \endgroup
  \if\relax\detokenize{#5}\relax
    \entrynobulletspace
  \else
    \entrybulletspace{#5}
  \fi
  \end{minipage}\par
  \entryspace
}

\newcommand{\entrywidedate}[5]{%
  \begin{minipage}[t]{\textwidth}
  \begingroup\setlength{\tabcolsep}{0pt}%
  \begin{tabular*}{\textwidth}{@{}p{\entrytitlewidth}@{\extracolsep{\fill}}r@{}}
    \textbf{#1} & #3 \\
    \textit{#2} & \mbox{#4} \\
  \end{tabular*}
  \endgroup
  \if\relax\detokenize{#5}\relax
    \entrynobulletspace
  \else
    \entrybulletspace{#5}
  \fi
  \end{minipage}\par
  \entryspace
}

\newcommand{\entrynarrowdate}[5]{%
  \begin{minipage}[t]{\textwidth}
  \begingroup\setlength{\tabcolsep}{0pt}%
  \begin{tabular*}{\textwidth}{@{}p{\entrytitlewidth}@{\extracolsep{\fill}}r@{}}
    \textbf{#1} & #3 \\
    \textit{#2} & \mbox{#4} \\
  \end{tabular*}
  \endgroup
  \if\relax\detokenize{#5}\relax
    \entrynobulletspace
  \else
    \entrybulletspace{#5}
  \fi
  \end{minipage}\par
  \entryspace
}

\newcommand{\awardcell}[2]{\textbf{#1}\par\textit{#2}}
\newcommand{\columntitle}[1]{\centering\textbf{\MakeUppercase{#1}}\\[-0.3em]{\color{accent}\rule{\linewidth}{0.4pt}}\arraybackslash}

\begin{document}

\begin{center}
  {\LARGE\bfseries\color{accent} Alessandro Dentico}\\[0.1em]
  {\color{accent}\rule{\textwidth}{0.4pt}}\\[0.35em]
  \small \href{mailto:alessandro.dentico01@universitadipavia.it}{\texttt{alessandro.dentico01@universitadipavia.it}} \textbullet\ \href{https://alessandro-dentico.github.io}{alessandro-dentico.github.io}\\[0.3em]
  \small \href{https://www.linkedin.com/in/alessandro-dentico}{linkedin.com/in/alessandro-dentico} \textbullet\ \href{https://orcid.org/0009-0000-0532-5080}{ORCID: 0009-0000-0532-5080}
\end{center}
\blockspace

\section*{Personal Profile}
First-year Ph.D. student in the joint programme between the University of Milan and the University of Pavia (coursework stage), specialising in health economics. My scholarship, financed by the Istituto Superiore di Sanit\`a (ISS), covers the use of statistical-econometric models for the analysis of population health and registry data; from the second year, my research will be carried out in cooperation with the ISS. So far, I have applied the Synthetic Control Method to evaluate the causal impact of Italy's Regional Recovery Plans on region-level healthcare staffing and patient out-migration.
\blockspace

\section*{Education}

\entrynarrowdate{University of Milan \& University of Pavia}{Joint Ph.D. Programme in Economics.}{Milan \& Pavia, Italy}{Oct. 2026 -- ongoing}{\begin{itemize}
  \item Thematic scholarship financed by the Istituto Superiore di Sanit\`a (ISS): \textit{``Use of Statistical--Econometric Models for the Analysis of Population Health Data and Registry Data''}.
  \item Year 1: coursework stage. | Year 2-4: research conducted in cooperation with the ISS, based in Rome.
\end{itemize}}

\entry{IUSS School for Advanced Studies}{Master of Advanced Studies in Social Sciences (``Master di II \mbox{livello}'').}{Pavia, Italy}{Sep. 2024 -- ongoing}{\begin{itemize}
  \item Attending additional interdisciplinary courses alongside the regular University curriculum.
  \item Relevant coursework: International Relations, History of Economic Thought, Comparative Analysis of Neoclassical Economic Models vs. Heterodox Models.
\end{itemize}}

\entrywidedate{University of Pavia}{M.Sc. in Economics, taught entirely in English.}{Pavia, Italy}{Sep. 2024 -- Jul. 2026}{\begin{itemize}
  \item Final grade: 110/110 \textit{summa cum laude} (equivalent to 100\% with honours).
  \item GPA: 29.55/30 (equivalent to 4.0).
  \item Thesis: \textit{Austerity, Workforce Cuts, and Patient Displacement: A Synthetic Control Analysis of Italy's Regional Recovery Plans}. | Supervisor: Prof. Cinzia Di Novi.
  \item Relevant coursework: Microeconometrics, Applied Econometrics, Policy Evaluation, Health Economics.
\end{itemize}}

\entrywidedate{Ludwig-Maximilians-Universit\"at M\"unchen (LMU)}{Visiting Student, guest of \textbf{Stiftung Maximilianeum}.}{Munich, Germany}{Apr. 2026 -- Jul. 2026}

\entry{University of Pavia}{B.Sc. in Economics.}{Pavia, Italy}{Sep. 2021 -- Oct. 2024}{\begin{itemize}
  \item Final grade: 110/110 \textit{summa cum laude} (equivalent to 100\% with honours).
  \item GPA: 29.41/30 (equivalent to 4.0).
\end{itemize}}

\clearpage

\section*{Additional Training}

\entry{European Economic Review (EER)}{3rd Summer School in Experimental and Behavioural Economics.}{Lisbon, Portugal}{Jun. 2026}{\begin{itemize}
  \item Covered behavioural game theory, social preferences, and biological factors in decision-making, with applications to incentives and policy design.
\end{itemize}}

\entry{National Model United Nations (NMUN)}{Model UN Conference New York 2026 (NMUN•NY).}{New York, USA}{Mar. 2026}{\begin{itemize}
  \item Completed a preparatory curriculum in public speaking and negotiation, then represented a Member State at the New York simulation.
\end{itemize}}

\entry{CEMFI (Centre for Monetary and Financial Studies)}{Summer School in New Developments in the Econometrics of Heterogeneous \mbox{Workers and Firms}.}{Madrid, Spain}{Sep. 2025}{\begin{itemize}
  \item Studied microeconometric methods to account for heterogeneity, including grouped fixed effects, clustering methods, and AKM decompositions.
\end{itemize}}

\entry{Sant'Anna School of Advanced Studies}{Summer School in Agent-Based Models in Economics.}{Pisa, Italy}{Jul. 2025}{\begin{itemize}
  \item Used the LSD platform to adapt an ABM to evaluate the impact of bank concentration.
\end{itemize}}

\entry{Sant'Anna School of Advanced Studies}{Summer School in Economics of Innovation and Technological Change.}{Pisa, Italy}{Jul. 2023}{\begin{itemize}
  \item Gained exposure to Neo-Schumpeterian Economics of Innovation.
\end{itemize}}

\section*{Professional Experience}

\entry{Collegio Ghislieri}{Rector's Office Assistant.}{Pavia, Italy}{Mar. 2024 -- Jul. 2025}{\begin{itemize}
  \item Managed and updated the college's digital information system.
  \item Developed digital solutions in Python to automate repetitive administrative tasks.
\end{itemize}}

\section*{Leadership \& Academic Involvement}

\entry{Collegio Ghislieri \& IUSS Pavia}{Organiser and Moderator of the Conference: Co-operation and Development in a Multipolar World.}{Pavia, Italy}{Mar. 2026}{}

\entry{Library Council for Social Sciences of the University of Pavia}{Member, elected as student representative.}{Pavia, Italy}{Jun. 2024 -- ongoing}{}

\entry{Council of the Department of Economics of the University of Pavia}{Member, elected as student representative.}{Pavia, Italy}{Jun. 2024 -- May 2026}{}

\clearpage

\section*{Scholarships, Prizes \& Funding}

\subsection*{Merit Scholarships}
{\rowcolors{1}{cellbg}{cellbg}\arrayrulecolor{white}\setlength{\arrayrulewidth}{2.4pt}\setlength{\tabcolsep}{3pt}%
\noindent\begin{tabularx}{\textwidth}{|>{\justifying\setlength{\parindent}{0pt}\arraybackslash}p{0.27\textwidth}|>{\justifying\setlength{\parindent}{0pt}\arraybackslash}X|>{\RaggedLeft\arraybackslash}p{0.10\textwidth}|>{\RaggedLeft\arraybackslash}p{0.11\textwidth}|}
\hline
\textbf{Collegio Ghislieri} & Full-ride scholarship awarded to those selected in a national selection, and renewed annually based on GPA maintenance, timely coursework completion, and 70 hours of internal additional training. & \textbf{Room\newline\& Board} & 2021--2026 \\
\hline
\textbf{IUSS School of Advanced Studies} & Annual merit stipend for the Master of Advanced Studies in Social Sciences, contingent on maintaining high GPA. & \textbf{\euro 5,000\newline annually} & 2024--2026 \\
\hline
\end{tabularx}}

\subsection*{Academic Prizes}
{\rowcolors{1}{cellbg}{cellbg}\arrayrulecolor{white}\setlength{\arrayrulewidth}{2.4pt}\setlength{\tabcolsep}{3pt}%
\noindent\begin{tabularx}{\textwidth}{|>{\justifying\setlength{\parindent}{0pt}\arraybackslash}p{0.27\textwidth}|>{\justifying\setlength{\parindent}{0pt}\arraybackslash}X|>{\RaggedLeft\arraybackslash}p{0.10\textwidth}|>{\RaggedLeft\arraybackslash}p{0.11\textwidth}|}
\hline
\textbf{University of Pavia}\par\textit{``Giovanni Manera'' Prize} & Awarded to the top-ranking Master's student across the Departments of Economics and Political Science to fund research abroad, based on academic record and thesis proposal quality. & \textbf{\euro 5,000} & Apr. 2026 \\
\hline
\textbf{University of Pavia}\par\textit{``Piero Bianco'' Prize} & Awarded for the highest GPA among all Bachelor's graduates in the Department of Economics. & \textbf{\euro 7,000} & Jun. 2025 \\
\hline
\end{tabularx}}

\subsection*{Research Mobility Grants}
{\rowcolors{1}{cellbg}{cellbg}\arrayrulecolor{white}\setlength{\arrayrulewidth}{2.4pt}\setlength{\tabcolsep}{3pt}%
\noindent\begin{tabularx}{\textwidth}{|>{\justifying\setlength{\parindent}{0pt}\arraybackslash}p{0.27\textwidth}|>{\justifying\setlength{\parindent}{0pt}\arraybackslash}X|>{\RaggedLeft\arraybackslash}p{0.10\textwidth}|>{\RaggedLeft\arraybackslash}p{0.11\textwidth}|}
\hline
\textbf{Stiftung Maximilianeum} & Mobility scholarship for a visiting studentship at LMU, Munich. & \textbf{Room\newline\& Board} & Apr.--Jul. 2026 \\
\hline
\textbf{IUSS School of Advanced Studies} & Funding to attend the Summer School at CEMFI, Madrid. & \textbf{\euro 1,500} & Sep. 2025 \\
\hline
\end{tabularx}}
\blockspace

\noindent{\rowcolors{1}{cellbg}{cellbg}\arrayrulecolor{white}\setlength{\arrayrulewidth}{2.4pt}\setlength{\tabcolsep}{3pt}%
\begin{tabularx}{\textwidth}{|X|X|}
\hline
\rowcolor{white}
\columntitle{IT Skills} & \columntitle{Language Skills} \\
\hline
\begin{minipage}[t]{\linewidth}
\setlength{\parskip}{0.45em}
\textbf{Proficient:} Stata, LaTeX, Notion.\par
\textbf{Good:} R, Office 365, Flourish, Canva.\par
\textbf{Basic:} Python.
\vspace{0.2em}
\end{minipage}
&
\begin{minipage}[t]{\linewidth}
\setlength{\parskip}{0.45em}
\textbf{Italian:} Native.\par
\textbf{English:} Proficient, C2 (Cambridge CAE: Grade A, taken in 2023; TOEFL iBT: 115/120, taken in 2025).\par
\textbf{French:} Intermediate.
\vspace{0.2em}
\end{minipage} \\
\hline
\end{tabularx}}
\vfill
{\footnotesize I hereby authorize the processing of the personal data contained in this CV for recruitment and selection purposes in accordance with EU Regulation 2016/679 (GDPR).}

\end{document}